\documentclass[11pt,a4paper]{article}

\usepackage[T1]{fontenc}
\usepackage[utf8]{inputenc}
\usepackage{lmodern}
\usepackage{microtype}
\usepackage[a4paper,margin=25mm,headheight=14pt]{geometry}
\usepackage{setspace}
\usepackage{enumitem}
\usepackage{booktabs}
\usepackage{tabularx}
\usepackage{longtable}
\usepackage{array}
\usepackage{xcolor}
\usepackage{amsmath,amssymb}
\usepackage{fancyhdr}
\usepackage{hyperref}

\definecolor{CQGblue}{RGB}{31,74,125}
\definecolor{DarkGreen}{RGB}{35,105,65}
\definecolor{DarkRed}{RGB}{145,38,38}
\definecolor{LightBlue}{RGB}{235,242,250}
\hypersetup{
  colorlinks=true,
  linkcolor=CQGblue,
  urlcolor=CQGblue,
  pdftitle={Editorial review of Paper I for Classical and Quantum Gravity},
  pdfauthor={Codex editorial review}
}

\setstretch{1.13}
\setlength{\parindent}{0pt}
\setlength{\parskip}{0.58em}
\setlist[itemize]{leftmargin=1.65em,itemsep=0.32em,topsep=0.35em}
\setlist[enumerate]{leftmargin=1.85em,itemsep=0.35em,topsep=0.4em}
\renewcommand{\arraystretch}{1.14}

\pagestyle{fancy}
\fancyhf{}
\lhead{Editorial review --- Paper I}
\rhead{CQG simulation}
\cfoot{\thepage}

\newcommand{\major}{\textcolor{DarkRed}{\textbf{Major}}}
\newcommand{\minor}{\textcolor{CQGblue}{\textbf{Minor}}}
\newcommand{\preserve}{\textcolor{DarkGreen}{\textbf{Preserve}}}
\newcommand{\manuscript}{\emph{Controlled-rail brachistochrones in non-stationary spacetimes: conformal symmetry, Kodama energy, and Vaidya dynamics}}

\begin{document}
\hypersetup{pageanchor=false}

\begin{titlepage}
  \centering
  \vspace*{2.0cm}
  {\Large\bfseries Editorial assessment for\\[0.4em] Classical and Quantum Gravity\par}
  \vspace{1.5cm}
  {\LARGE\bfseries \manuscript\par}
  \vspace{1.1cm}
  {\large Review of the 38-page file \texttt{paper1-11.pdf}\par}
  \vspace{1.8cm}
  \begin{minipage}{0.84\textwidth}
    \colorbox{LightBlue}{\parbox{0.94\textwidth}{
      \textbf{Simulated editorial decision: major editorial revision before external review.}
      The manuscript is within CQG's scope and contains a potentially publishable,
      distinctive contribution.  In its present form, however, the scientific story is
      obscured by an overextended technical layer, incomplete positioning against the direct
      relativistic-brachistochrone literature, and several visible inconsistencies between
      carefully qualified body text and stronger captions or summary claims.
    }}
  \end{minipage}
  \vfill
  {\large Prepared as an editorial review, not as a substitute mathematical referee report\par}
  \vspace{0.4cm}
  {\large 3 August 2026\par}
  \vspace{1.3cm}
  {\small Codex-origin file; no manuscript source was changed during this assessment.\par}
\end{titlepage}
\hypersetup{pageanchor=true}
\pagenumbering{roman}

\tableofcontents
\clearpage
\pagenumbering{arabic}

\section{Mandate, material reviewed, and limits of this report}

This report follows the supplied editorial-review brief.  I read the complete manuscript,
including all three appendices, the data-availability statement, acknowledgements and the
69-entry reference list.  I also rendered and inspected all 38 pages so that comments on
figures, tables, page flow and typography are based on the actual PDF rather than only on
extracted text.  The repository source was consulted to resolve cross-references, citation
usage and the relationship to the companion paper.  Current CQG author guidance and IOP's
generative-AI policy were checked, as was directly relevant relativistic-brachistochrone
literature.

The task is editorial.  I therefore assess scope, novelty presentation, readability,
scientific narrative, claim calibration, figures, notation, bibliography, research-integrity
statements and the probability of passing an initial editorial screen.  I do not rederive the
Hamiltonians, certify the special-function reductions, or replace a specialist referee's
judgement on the validity of Theorems I.1--I.5.  Where a mathematical statement creates an
editorial problem, I identify the mismatch or ambiguity and recommend safer wording; I do not
declare a new mathematical result.

The authoritative object for this review is the supplied \path{paper1-11.pdf}.  The local
\path{paper1/paper1.tex} source appears slightly ahead of that PDF in at least one
cross-reference, so final production must include a clean recompilation and PDF-to-source
comparison.

\subsection*{Severity vocabulary}

\begin{description}[leftmargin=1.8cm,style=nextline]
  \item[\major] An issue likely to affect desk review, the perceived novelty, the central
  narrative, or the reliability of the claims.  It should be resolved before submission.
  \item[\minor] A clarity, consistency, presentation or metadata issue that should be fixed but
  does not determine publishability by itself.
  \item[\preserve] A feature that is already effective and should not be lost during revision.
\end{description}

\section{Executive summary}

The manuscript asks a clear and worthwhile question: how should a relativistic
brachistochrone be formulated when the spacetime itself evolves and no timelike Killing energy
is available?  Its answer is to treat the rail as an active control constraint,
$-u\!\cdot\!W=\hat E$, and to select $W$ through a Killing $\rightarrow$
conformal-Killing $\rightarrow$ Kodama hierarchy.  It then develops a Pontryagin formulation,
an existence/normality statement on a regular timelike-selector domain, a conditional
Hamilton--Jacobi verification criterion, an FLRW base case and a detailed ingoing-Vaidya
application.  The Vaidya material includes exact rail kinematics, arrival/proper-time branch
structure, plunge and timing phenomena, and a first-order adiabatic correction whose complete
on-shell plus off-shell form is checked against the true non-autonomous flow by the expected
quadratic residual scaling.

This is recognisably a CQG paper.  It concerns classical general relativity, dynamical black-hole
geometry, constrained Hamiltonian systems and mathematical methods, all explicitly within the
journal's remit.  The Kodama-based controlled invariant is a natural bridge between dynamical
spacetime geometry and optimal control.  The manuscript also displays unusually good habits in
several places: it distinguishes the conserved Pontryagin costate from mechanical angular
momentum; states that the ordinary Hamiltonian is not an interior first integral; separates
PMP extremality from global HJB certification; labels the negative-rate ingoing-Vaidya branch as
a formal continuation rather than physical evaporation; archives code and data; and divides
the conclusion into proved, conditional/numerical and open statements.

The principal editorial problem is that two papers still coexist inside Paper I.  One is a
broadly readable conceptual and physical article about controlled relativistic travel in FLRW
and Vaidya.  The other is a specialist mathematical article about genus-two Abelian integrals,
third-kind horizon letters, iterated integrals, elliptic degeneration, theta-nome series and
conjectural transcendence.  Both may be valuable, but the second currently overwhelms the first
in pages 19--22 and Appendices B--C.  The main physical result is consequently buried, and CQG's
broad gravitational readership is asked to traverse the full special-function taxonomy before
receiving a compact statement of what has been proved and what it means.

A second major problem is literature positioning.  The introduction cites Perlick and general
Randers/Zermelo sources but omits the direct sequence of relativistic-brachistochrone papers by
Giannoni, Piccione, Verderesi and Tausk on sub-Riemannian formulations, existence, travel-time
Morse theory and arrival-time brachistochrones.  The sentence that the existence question and
the appropriate replacement invariant ``have not been addressed'' is therefore too broad.
The novelty remains defensible, but it must be stated more precisely: the new element is the
active maintenance of a selected rail charge in a genuinely non-autonomous background, the
Kodama selector, the explicit free-arrival Pontryagin system and its FLRW/Vaidya realisation.

The third major problem is internal hierarchy.  Theorem I.5, which supplies the complete
first-order extended-Hamiltonian correction needed to close the earlier off-shell objection,
appears only on page 32 in Appendix C.  Lemma I.6 appears before it on page 29, so results are
encountered in the order I.4, I.6, I.5.  At the same time, captions for figures 3 and A2 use
``minimizes'' after only a restricted perturbation check, although section 2 carefully explains
that generic global optimality is not certified.  The final PDF also contains a stale-looking
equation reference in the caption of figure 10.  These are exactly the kinds of inconsistencies
that reduce an editor's confidence even when the underlying derivation is sound.

Finally, the paper still leaks substantial Paper-II material.  The notation table contains the
Thakurta factor, Kerr spin and conformally normalised momenta; figure 1 includes a Thakurta--Kerr
panel; ergosphere caveats recur; and many unused Kerr/cosmological references are printed because
the source ends with \verb|\nocite{*}|.  There are 69 listed references but only 40 explicitly
cited keys in the Paper-I text.  This makes the split appear incomplete.

The recommended revision is therefore not a new scientific programme.  It is a decisive
editorial reassembly of results that are largely already present:

\begin{enumerate}
  \item rewrite the novelty paragraph around the precise controlled non-stationary contribution
  and add the direct brachistochrone literature;
  \item bring the statement and physical meaning of Theorem I.5 into the main text;
  \item reduce the special-function body section to the result, class of representation and
  validation, moving detailed construction to appendices or supplementary material;
  \item remove most Paper-II-specific notation, imagery and uncited bibliography;
  \item make every claim of minimisation consistent with the manuscript's own PMP/HJB hierarchy;
  \item reorder floats so that the Conclusions are uninterrupted and rebuild the cramped
  reproducibility table.
\end{enumerate}

After those changes, I would expect the paper to be a plausible candidate for external review.
In its present form I would return it for substantial editorial revision rather than reject it
for scope or lack of substance.

\section{Simulated editorial decision}

\subsection{Decision: major editorial revision before review}

If acting as an Associate Editor at initial screening, I would not desk-reject the manuscript on
scope.  The topic is appropriate, the paper has a coherent central idea, and the Vaidya result
could interest both gravitational theorists and mathematical relativists.  I would, however,
ask for a substantially revised manuscript before recruiting referees.

The reason is practical as well as scientific.  A specialist referee must be able to identify the
claim being reviewed.  Here the core claim changes scale several times: a well-posed controlled
problem, local PMP extremals, conditional global HJB verification, frozen phenomenology,
first-order non-autonomous closure, then detailed claims about transcendental classes.  The
manuscript contains many correct caveats, but they are distributed across captions, appendices
and parenthetical remarks.  The editorial burden of reconstructing the hierarchy is too high.

\subsection{Desk-review risk profile}

\begin{tabularx}{\textwidth}{@{}p{0.26\textwidth}X p{0.20\textwidth}@{}}
\toprule
Criterion & Assessment & Risk now \\
\midrule
Journal scope & Strong fit: classical GR, dynamical black holes, spacetime geometry and
Hamiltonian control. & Low \\
Original contribution & Potentially strong, but novelty is framed too broadly and not compared
with direct predecessors. & Moderate--high \\
Broad-reader motivation & Good opening question and physical applications; operational meaning
of the rail appears too late. & Moderate \\
Scientific claim hierarchy & Careful in section 2 and the conclusion, but not uniform in abstract,
captions and appendix placement. & High \\
Length and focus & No hard CQG limit is assumed here, but the content-to-length ratio and two-paper
feel are problematic. & High \\
Presentation quality & Generally clean typesetting; figures and tables need a journal-facing
redesign and better float order. & Moderate \\
Reproducibility/integrity & Strong archive and disclosure posture; exact snapshot and bibliography
curation still need checking. & Low--moderate \\
\bottomrule
\end{tabularx}

The likely current outcome depends on the handling editor.  A mathematically inclined editor
might send it out because of the evident work and the careful qualifications.  A broad-scope
editor might return or reject it because the novelty comparison is incomplete and the special
function layer obscures the gravitational advance.  The proposed revision is aimed at removing
that avoidable variance.

\section{Principal strengths to preserve}

\subsection{A distinctive organising idea}

The selector hierarchy is memorable.  Killing, conformal Killing and Kodama structures are not
presented as unrelated tricks but as progressively weaker geometric data used to define the
controlled charge.  This gives the work a conceptual spine that can support a series of papers.
The hierarchy should remain visible in the abstract and introduction, even if Paper-II-specific
detail is reduced.

\subsection{An explicit experimental protocol}

The paper now states what is fixed and what is free: a launch event, a spatial target and a rail
charge are prescribed, while the arrival clock is free.  Equation (3) is correctly identified as
a free-final-clock transversality condition rather than an identity imposed throughout a
non-autonomous arc.  This distinction was evidently hard won and should not be diluted during
copy-editing.

\subsection{Honest separation of necessary and sufficient optimality statements}

Section 2 says plainly that PMP supplies necessary conditions, whereas an HJB value function is a
sufficient global certificate.  It also discusses conjugate points, Maxwell sets and cut loci,
and admits that a global value function is not constructed for generic non-stationary Vaidya or
Kerr trajectories.  This is excellent scholarly hygiene.  The required revision is to make the
rest of the paper obey the same vocabulary.

\subsection{The Vaidya control interpretation}

Using the Kodama vector as the unnormalised selector is physically and geometrically natural in
spherical symmetry.  The manuscript does more than write down a formal charge: it distinguishes
controlled conservation from Noether conservation and identifies the proper acceleration/power
required to maintain the rail.  This makes the construction falsifiable as a protocol rather
than merely a relabelling of geodesics.

\subsection{Correction of the off-shell issue}

The complete first-order response is now based on the extended Hamiltonian in the original
canonical variables, includes the interior momentum-shell shift, and introduces the anchored
source $S_{\mathcal D}$.  For Vaidya it collapses to the useful boundary form
$S_{\mathcal D}=[r p_r]-\lambda$.  The log--log diagnostic in figure 10 distinguishes the
$O(\varepsilon)$ residual of the on-shell-only term from the approximately
$O(\varepsilon^2)$ residual of the complete first-order term.  Editorially, this is a major
strength: it converts a previous objection into a clear narrative of correction and verification.

\subsection{Result-tiered conclusion and open-problem discipline}

The ``Proved / Conditional and numerical / Outlook'' division is worth preserving nearly
verbatim in spirit.  The manuscript also resists presenting the negative-rate ingoing chart as a
physical evaporating black hole, and states that the true outgoing-Vaidya boundary-value problem
remains open.  That distinction protects the paper from a serious physical overclaim.

\subsection{Reproducibility and AI disclosure}

The data/code DOI, script map, reported residuals and explicit acknowledgement of AI-assisted
symbolic, numerical and editorial work are positive features.  Current IOP guidance permits such
uses when disclosed and critically verified by the author.  The statement already identifies
models and purposes and assigns responsibility to the author; it should be retained, with only
minor tightening.

\section{Principal weaknesses and why they matter}

\subsection{The manuscript has two centres of gravity}

The controlled-rail/Kodama story and the higher-genus closure story compete rather than support
each other.  A broad CQG reader needs only enough special-function information to know that the
first-order response is explicit, evaluable, and structurally more complicated because of the
horizon logarithm.  Instead, the main text develops local uniformisers, residue divisors,
theta-ratio representations, elliptic degeneration data, rank-five numerical evidence and
conjectural irreducibility.  This is specialist content, and some of it concerns a double root at
negative radius explicitly declared nonphysical.

The cost is not merely length.  It makes the physical advance look like motivation for a special-
functions paper rather than the outcome of a gravitational-control paper.  It also forces the
reader to remember whether each formula belongs to the on-shell component, the complete response,
the algebraic degeneration family, the physical branch or a frozen advanced/retarded comparison.

\subsection{The novelty claim is vulnerable}

The phrase that the well-posedness question and appropriate invariant ``have not been addressed''
is broader than the bibliography supports.  Earlier work studied relativistic brachistochrones
with prescribed energy relative to observer fields, sub-Riemannian reductions, existence and
multiplicity, and distinctions between travel and arrival time.  General Fermat principles also
do not require stationarity in their most abstract spacetime formulation.

None of this eliminates the manuscript's novelty.  It changes the comparison.  The paper should
claim the controlled non-autonomous charge, canonical Kodama choice, explicit dynamic
Hamiltonians and first-order response, not the invention of relativistic arrival-time
variational theory.  Failure to make that distinction could trigger a desk rejection or an
adverse first referee report even if all calculations are correct.

\subsection{The central theorem is buried and out of order}

The complete first-order response is central because it justifies the headline statement that the
adiabatic correction has been closed including the off-shell term.  Yet Theorem I.5 first appears
on page 32.  Before it, the reader encounters Lemma I.6 on page 29.  This ordering suggests that
the appendices were assembled from separate development streams rather than designed around the
final logic.

The theorem statement belongs in the main text.  Its derivation may remain in an appendix.  The
body should then cite an appendix for the proof, not ask the reader to discover the theorem after
the conclusion and reproducibility appendix.

\subsection{Qualification is not consistent across text and captions}

Section 2 correctly says that a generic PMP extremal is not automatically a global minimiser.
Nevertheless, figure 3 says the straight line ``minimizes both'' functionals after a displayed
one-parameter perturbation, and figure A2 says each branch ``minimizes its own travel time.''  A
restricted perturbation family establishes a useful sanity check, not global minimality over the
full admissible class.

The conclusion's opening ``well-posed'' statement also needs the compact-domain/timelike-selector
qualification present in Theorem I.2.  These are editorial inconsistencies with mathematical
consequences.  They can be fixed by disciplined vocabulary, without changing equations.

\subsection{Paper-I/Paper-II leakage remains visible}

Table 1 devotes rows to $A(\eta)$, Kerr spin, $E_{\rm eff}$, $J_{\rm eff}$ and $P_r$; figure 1
contains a Thakurta--Kerr panel; and pages 4, 9, 10 and 32 revisit ergosphere or conformal-momentum
issues.  A small amount of forward reference helps establish the series, but the current level
suggests the split was performed mechanically.

The bibliography confirms the impression.  Because \verb|\nocite{*}| is used with the shared
bibliography, the PDF prints 69 entries although only 40 citation keys occur in the text.  Many
uncited entries concern Kerr, Thakurta, McVittie or mathematical tools not invoked in Paper I.

\subsection{Figures dominate the page flow without a single editorial hierarchy}

The plots are useful for research verification, but not all are main-story figures.  Full-page
panels interrupt text; verbose plot titles duplicate captions; and conclusion paragraphs are
separated by figures 8--10.  Figure 7 spends more ink limiting its own interpretation than
communicating a publishable result.  Figure 1 is too small to read comfortably, while figures 2
and 3 are larger than their narrative weight warrants.  The current sequence resembles a lab
notebook exported into a paper.

\section{Section-by-section editorial comments}

\subsection{Title and first page}

The current title is accurate but long.  It combines a new term (controlled rail), a general
geometric mechanism (conformal symmetry), a spherical dynamical mechanism (Kodama energy) and a
specific application (Vaidya).  A reader cannot tell which is primary.  If the revised paper is
application-led, I recommend \emph{Time-optimal controlled worldlines in FLRW and Vaidya
spacetimes}.  If the selector hierarchy remains central, I recommend \emph{Controlled-rail
brachistochrones in dynamical spacetimes: from conformal symmetry to Kodama flow}.

The first page has substantial lower whitespace, but this is not itself an editorial defect: the
IOP submission guide allows flexible formatting and production will re-typeset accepted work.
Do not spend revision effort filling that space.  Spend it clarifying the title and abstract.

The PDF has no keywords.  IOP collects them at submission, so this is not a formal omission in the
manuscript, but the author should prepare a consistent set: relativistic brachistochrone; optimal
control in general relativity; Kodama vector; Vaidya spacetime; Randers--Finsler geometry;
non-autonomous Hamiltonian systems; dynamical horizons.

\subsection{Abstract}

The abstract is within the usual 300-word limit and is self-contained.  Its opening question is
effective.  The problem is density and calibration.  Within a single paragraph it introduces the
rail, selector, three-level hierarchy, PMP, existence, normality, HJB, FLRW, Vaidya, special
functions, true-flow verification and the companion paper.

Three phrases should be changed:

\begin{itemize}
  \item ``legitimate Pontryagin problem---with existence, normality, and a ... verification
  criterion'' should not imply that generic global optimality has been proved;
  \item ``solve two ... cases in closed form'' should be narrowed to the derivation of extremal
  equations and closed representations at the levels actually established;
  \item ``cleanly separates'' is evaluative; state the concrete separation between spatial-gradient
  and mass-flow terms.
\end{itemize}

An alternative abstract is provided in section~\ref{sec:modeltext}.

\subsection{Section 1: Introduction}

The introduction is commendably brief and its four-paragraph structure is sound.  Preserve the
opening stationary/dynamical contrast and the explicit paper map.  Revise the second paragraph,
which presently turns a legitimate research gap into an absolute historical claim.

The reader also needs the operational meaning of the rail before the formalism.  Add one sentence
explaining that an ideal actuator supplies the proper acceleration needed to hold the chosen
observer-relative energy fixed, and that the cost being minimised is elapsed clock time, not
actuation energy.  This prevents a predictable objection: if arbitrary thrust is allowed, why is
the problem not trivial?  The answer lies in the fixed indicatrix/charge constraint, but it should
be visible immediately.

The last paragraph gives too much promotional space to the companion paper.  One sentence is
sufficient.  Paper I should end its introduction by foregrounding its own strongest result: the
complete first-order Vaidya response and its true-flow test.

\subsection{Section 2.1: Rail constraint and indicatrix}

The sequence rail constraint $\rightarrow$ ellipse $\rightarrow$ domain $\rightarrow$ protocol
$\rightarrow$ PMP is logical.  The distinction between the boundary of a convex region and a
convex set is careful and should be preserved.  So should the statement that no relaxed control is
needed because the maximising direction is unique on the strictly convex oval.

The section is nonetheless too dense for the first encounter.  Table 1 arrives before the reader
has seen either application and includes Paper-II notation.  Reduce it to $u,W,\hat E,J$,
$L_{\rm mech}$, the clocks and the adiabatic parameter.  Introduce Vaidya-specific symbols at the
start of section 4 and leave conformal Kerr symbols to Paper II.

The ``Domain of the indicatrix'' paragraph simultaneously discusses freezing, horizons,
ergospheres, analytic continuations and Doran/ZAMO alternatives.  Only the general timelike-selector
and nondegeneracy assumptions are needed here.  Move the rotating caveat to the companion paper.

The phrase ``fixed endpoints'' must be standardised.  The protocol is a fixed initial event plus a
fixed spatial target/terminal worldline at free arrival time, not two fixed spacetime events.  When
Theorem I.2 or figure A2 uses ``fixed endpoints'', specify ``fixed spatial endpoints'' or the exact
terminal manifold so it cannot be confused with the companion two-event problem.

\subsection{Section 2.2: Existence, normality and global optimality}

This section contains essential corrections and is one of the strongest parts of the manuscript.
The direct-method assumptions are stated; the freezing/horizon exclusion is visible; normality is
argued from the positive support function; and the HJB result is identified as a verification
theorem rather than an automatically available solution.

Editorially, the proof detail can move to an appendix.  The main text needs a boxed theorem with
the domain and conclusion, followed by one paragraph stating what is and is not globally
certified.  The extended discussion of viscosity solutions, conjugate points, Maxwell sets and
cut loci is useful for mathematical readers but slows the physical narrative before FLRW.  A
compact summary in the body and a fuller appendix would preserve rigour while improving flow.

Avoid the sentence ``a candidate extremal saturates it'' unless the saturation relation and
regularity assumptions are immediately clear.  Broad readers may interpret HJB terminology as a
completed global solution.  The most important sentence is already present: no global value
function is constructed for generic non-stationary trajectories.  Bring it forward.

\subsection{Section 2.3: Selector hierarchy and Perlick limit}

This is conceptually valuable and should remain in the main text.  Recovery of Perlick's fixed-
energy stationary optics is an important consistency check.  However, the derivation can be
shortened after the equality of the Finsler functions is stated.  The hierarchy should be
introduced as a convention for the controlled problem, not as the only possible physically
canonical choice in all spacetimes.  The current final ``Convention'' item does this honestly;
retain it.

The conformal-Killing level should be described only to the extent needed for FLRW.  References to
constant-$A$ Thakurta--Kerr and detailed conformal normalisation belong in the companion article.

\subsection{Section 3: FLRW}

The FLRW result is clear and appropriately simple: the position-independent optical index makes
the straight comoving line the shared spatial path, with freezing but no branch splitting.  This
serves as an effective base case for the Vaidya contrast.

Two points need care.  First, equation (10) uses $\arg\min$ notation for all three clocks.  If the
underlying argument proves the equality over the full admissible class, state its assumptions.  If
figure 3 is merely a restricted perturbation check, do not use that figure as the proof.  Second,
the phrase ``variational check: ... minimizes BOTH'' inside the plot and its caption overstates what
a single family of sinusoidal perturbations establishes.  Use ``local perturbation check''.

Figures 2 and 3 can be combined into a two-panel figure.  The section would then occupy roughly one
page and function as the intended bridge rather than a two-page visual interlude.

\subsection{Section 4.1: Vaidya invariant and Hamiltonians}

The transition to Vaidya is effective.  The Kodama vector and Misner--Sharp mass are introduced with
standard references, and the controlled/Noether distinction is explicit.  The branch Hamiltonians
and exact map between the axial costate and mechanical angular momentum are essential; preserve
them.

The paragraph on on-shell drift versus off-shell shell departure is particularly important because
it prepares the later adiabatic correction.  It should be connected more directly to Theorem I.5:
state here that the full first-order response must include the induced costate displacement and
give a forward reference to the theorem statement in the main text.

Figure 4 should be assessed after the final layout at actual column size.  The information is
useful, but the typography should match the later Vaidya figures and avoid a large explanatory
headline inside the image.

\subsection{Sections 4.2--4.3: Phenomenology and frozen-clock comparison}

Penetration, arrival-epoch sensitivity and regular integration through periapsis are physically
readable and likely to attract CQG readers.  These results should be narrated before the detailed
algebraic machinery.  A compact summary table could distinguish: exact identity; frozen
comparison; full non-autonomous numerical result; and open outgoing problem.

The negative-rate branch is now correctly labelled as formal.  Even so, figures 7--9 grant it too
much prominence.  Figure 7 in particular uses a full page to state that a sign change occurs only
after entering an unphysical negative-mass region and that the outgoing problem is not solved.
That conclusion can be conveyed in three sentences.  Move the plot to supplementary material or
retain only one panel as a robustness check.

The term ``absence of inversion'' should always carry ``in the frozen-clock comparison''.  The
current subsection title does this; captions and conclusion should use the same qualifier.  Do not
let a reader mistake the result for a theorem about physical evaporation.

\subsection{Section 4.4: Closed first-order correction}

This section contains both the paper's most important technical result and its greatest editorial
obstacle.  The beginning is good: it identifies the frozen Schwarzschild family, the genus-two
sextic, the mass derivative, the two clocks and the horizon logarithm.  The paragraph that labels
equation (22) as the on-shell component and then explains the missing off-shell costate term is
essential.

The next several pages go too far for the main text.  A journal reader does not need all of the
following before reaching the conclusion: the complete divisor of the horizon differential,
local uniformisers, a sample negative-radius degeneration, Weierstrass sigma/zeta formulae, a
rank-five numerical basis claim, Kronecker--Eisenstein series, shuffle reductions and a detailed
advanced/retarded decomposition.  These are suitable for appendices, a supplement, or a separate
mathematical methods paper.

The main section should instead contain five objects:

\begin{enumerate}
  \item a precise definition of the frozen zeroth-order orbit;
  \item a statement of the complete first-order theorem, including the off-shell source;
  \item the Vaidya boundary reduction $S_{\mathcal D}=[r p_r]-\lambda$;
  \item a compact statement that the result is evaluable in a finite Abelian/iterated-integral
  basis, with the horizon logarithm adding a weight-two letter;
  \item figure 10 and enough numerical detail to interpret the slopes.
\end{enumerate}

Use ``closed representation'' rather than ``closed form'' whenever evaluation still requires
constructing a period matrix and using the archived theta routine.  The current paragraph ``What
is reconstructible from the text alone'' is admirably honest and should be retained, preferably
near the first use of ``closed'' rather than deep in Appendix C.

Claims about irreducibility are correctly marked conjectural in several places.  Enforce that
qualifier everywhere, including the abstract-level description if the transcendence class is
mentioned at all.  The numerical rank-five observation should not appear in the main narrative;
it is a secondary structural experiment.

\subsection{Section 5: Conclusions}

Preserve the three-tier structure but move every figure before the section.  In the supplied PDF,
``Proved'' appears below figure 8 on page 23, ``Conditional and numerical'' below figure 9 on page
24, and ``Outlook'' below figure 10 on page 25.  This fragments the final argument and makes the
reader interpret plot captions between conclusion paragraphs.

Qualify ``well-posed'' by the compact timelike-selector domain and stated endpoint protocol.
Replace ``the first-order correction closes in a horizon dilogarithm'' with a formulation that
distinguishes the complete response from the special-function representation of its components.
End on the Vaidya conclusion; the companion rotating work should receive only a final sentence.

\subsection{Appendix A: Reproducibility and consistency}

The exact polynomial identity and script map are valuable.  They allow a referee to distinguish
symbolic derivation from fitting.  Keep this appendix, but streamline it.

Table A2 is visibly broken at the bottom: the extended-Hamiltonian row, script name, slope and hash
are cramped and partly collide.  Rebuild it with one result per row, a fixed-width or
\texttt{tabularx} description column, and separate columns for script, test quantity and achieved
residual.  Put paths and hashes in a smaller monospaced font.  Verify that the script hash matches
the exact Zenodo version cited by the paper.

Figure A1 is a useful stationary-limit check.  Figure A2 again uses ``minimizes'' too strongly;
change it to a restricted perturbation check.  If space is needed, move A2 to supplementary
material because its message duplicates figure 3 at a different geometry.

\subsection{Appendix B: Genus-degeneration loci}

This appendix is mathematically rich but editorially overdeveloped relative to its physical role.
It explicitly says that the Vaidya double root is at negative radius and is not a physical external
separatrix.  The manuscript should therefore avoid calling this simply ``the separatrix'' in
nearby prose or filenames.  Use $J_{\rm deg}$ and ``algebraic genus-degeneration locus''
consistently; reserve ``separatrix'' for a dynamical boundary when one exists.

The appendix would work better as supplementary mathematical material.  If it remains in the
article, begin with a short statement of why a CQG reader needs it: it proves that the generic
genus-two representation has a controlled genus-one limit and supplies an independent check.
Then move implementation-level residue formulae and numerical-basis compression to the archive.

Lemma I.6 appears before Theorem I.5.  Either reorder Appendices B and C, move the theorem statement
to the body, or renumber the results in reading order.  The first option is easiest structurally;
the second is best narratively.

\subsection{Appendix C: Complete extended-Hamiltonian correction}

The content of Appendix C is central, not optional.  Its theorem statement should be promoted to
section 4.4 or to the end of section 2.  The appendix can then be titled ``Derivation of the
complete first-order response'' and contain the canonical calculation, the drift generator, the
boundary source, branch orientation and separatrix limit.

The notation is demanding because $s$ is an evolution clock, $\lambda$ is elapsed from the
free-arrival endpoint, and $m(v)=m_0(1+\varepsilon s+\cdots)$ uses both $v$ and $s$.  Add a one-line
orientation diagram or a small table.  State the direction of integration once and reuse it.

The paragraph on Thakurta--Kerr conformal momentum normalisation is more detailed than Paper I
needs.  Replace it with a sentence saying the companion application adds a dilation term when
normalised momenta are time-dependent.  This preserves the logical bridge without reproducing
Paper II.

\subsection{Back matter}

The data-availability statement is strong.  Add the exact archive version or record version date
if the DOI resolves to multiple versions.  The GitHub link may remain as development history, but
the archived DOI must be sufficient to reproduce the submitted PDF.

The AI acknowledgement is acceptable in substance under current IOP guidance.  A concise version
could read: ``OpenAI GPT-5.6 (Sol) and Anthropic Claude (Opus 4.8) assisted with symbolic-algebra
checks, numerical-code development, and language/structural editing.  The author critically
reviewed and independently verified all outputs and takes responsibility for the work.''  Retain
the software citations separately.

The reference list must be regenerated without \verb|\nocite{*}| and checked against the actual
text.  Companion-paper material should not be imported wholesale.

\section{Literature positioning and novelty repair}

\subsection{The missing direct lineage}

The paper currently moves from Perlick (1991) to broad Randers/Zermelo and optimal-control sources.
This omits work that is terminologically and variationally much closer to the present article.
At minimum, the revised introduction should discuss:

\begin{itemize}
  \item Giannoni, Piccione and Verderesi (1997), on a sub-Riemannian formulation of the relativistic
  brachistochrone;
  \item Giannoni and Piccione (1998), on existence in stationary spacetimes;
  \item Giannoni, Piccione and Tausk (2002), on travel-time brachistochrones, observer-field
  conservation laws, first/second variation and Morse theory;
  \item Giannoni and Piccione (2002), on arrival-time brachistochrones in a general relativistic
  setting and the distinction from travel-time problems.
\end{itemize}

These sources do not obviously contain the controlled Kodama construction.  Their relevance is
that they occupy the very conceptual territory currently described as unaddressed: observer-field
constraints, prescribed energy, arrival/travel time, endpoints and existence.  Citing them allows
the author to articulate the real advance with much greater credibility.

\subsection{General Fermat principles are broader than stationary Randers reduction}

The manuscript is correct that the stationary spatial reduction to a time-independent Randers or
Jacobi geometry depends on symmetry.  It should not imply that general-relativistic Fermat
principles themselves exist only in stationary spacetimes.  Arrival-time characterisations of
null geodesics can be formulated on arbitrary spacetimes, and there are timelike extensions with
prescribed proper time.  The controlled massive rail problem is different, but the distinction
must be explicit.

A useful 2026 survey by Caponio and Javaloyes synthesises these formulations and can anchor the
terminology.  It may be cited as a survey rather than a priority claim.

\subsection{Very recent adjacent work}

The July 2026 preprint \emph{The brachistochrone problem for a constant velocity traveler in static
and spherically symmetric spacetimes} studies a fixed local-speed protocol in stationary spherical
geometries.  Because the supplied manuscript is dated August 2026, this work is recent but already
visible.  It should be considered in the final literature search.  A short contrast would be
enough: fixed local speed in a static observer frame is not the same admissible set as a fixed
Killing/Kodama rail charge in a dynamical spacetime.

Another recent paper on brachistochrone-ruled timelike surfaces uses stationary Finsler/Jacobi
reductions.  It is adjacent rather than essential.  Cite it only if the revised introduction
discusses current stationary applications; do not expand the bibliography merely for currency.

\subsection{Proposed novelty statement}

The following is both narrower and stronger than the current absolute claim:

\begin{quote}
Previous relativistic brachistochrone theories exploit stationarity, prescribed observer-field
energy, or geodesic arrival-time principles.  We consider instead an actively controlled timelike
worldline in a genuinely evolving geometry.  The control maintains a fixed charge relative to a
geometrically selected field, using the Kodama vector in spherical symmetry, and leads to a
non-autonomous free-arrival Pontryagin problem.  We derive this problem explicitly in FLRW and
ingoing Vaidya and obtain its complete first-order adiabatic response.
\end{quote}

This statement invites a clean comparison rather than a historical challenge.

\section{Recommended scientific narrative}

The revised paper should be organised around a sequence of questions a CQG reader naturally asks.

\subsection*{Question 1: what is the physical protocol?}

Define the rail in operational terms.  An actuator maintains $-u\cdot W=\hat E$; the control is
the direction on the resulting velocity indicatrix; the objective is a selected arrival or proper
clock; thrust expenditure is not the cost.  State the endpoint data and domain.

\subsection*{Question 2: why is the construction mathematically legitimate?}

State the PMP reduction, existence/normality domain and free-arrival transversality.  Give the HJB
criterion and immediately state that generic global certificates are not constructed.  Put proof
detail in an appendix.

\subsection*{Question 3: what geometric field supplies the charge?}

Present the Killing/CKV/Kodama hierarchy and the Perlick stationary limit.  Use FLRW as the
one-paragraph CKV example and Vaidya as the main Kodama example.  Mention Paper II only once.

\subsection*{Question 4: what does time dependence change physically?}

Lead with the Vaidya results: the radial wind offsets the indicatrix; arrival and proper-time
branches split; penetration and periapsis depend on the epoch; bounce integration remains regular.
Clearly label exact, frozen and numerical statements.

\subsection*{Question 5: is the first-order response actually the true non-autonomous one?}

State Theorem I.5 in the body.  Explain in plain language that differentiating the frozen shell is
not enough because the interior trajectory is off that shell; the anchored costate source supplies
the missing term.  Show the universal Vaidya boundary form and figure 10.  This is the natural
climax of the paper.

\subsection*{Question 6: how explicit is the answer?}

Give a compact hierarchy:

\begin{itemize}
  \item algebraic coefficient reduction: exact;
  \item weight-one Abelian/Kleinian terms: explicit;
  \item weight-two iterated/horizon terms: closed representation with convergent numerical
  evaluation;
  \item irreducibility/minimal basis: conjectural or numerical evidence only.
\end{itemize}

Then refer to appendices/supplement for construction.  End with limitations and the physical
outlook.

\subsection{Concrete revised outline}

\begin{longtable}{@{}p{0.08\textwidth}p{0.32\textwidth}p{0.52\textwidth}@{}}
\toprule
Part & Proposed heading & Purpose \\
\midrule
\endhead
1 & Introduction & Problem, direct literature, operational rail, precise contributions. \\
2 & Controlled-rail optimal control & Protocol, indicatrix, selector hierarchy, compact theorem
statements and qualifications. \\
3 & FLRW base case & One equation, one combined figure, branch degeneracy and freezing. \\
4 & Kodama-controlled Vaidya brachistochrones & Invariant, Hamiltonians, costate map and physical
phenomenology. \\
5 & Complete first-order response & Theorem I.5, boundary source, special-function hierarchy and
true-flow validation. \\
6 & Discussion and limitations & Meaning of control cost, globality, outgoing Vaidya, domain and
turning-point limitations. \\
7 & Conclusions & Uninterrupted, result-tiered summary. \\
A & Proofs of the general control statements & Direct method, normality, HJB and cut-locus detail. \\
B & Reproducibility & Polynomial reduction, script map, residual specifications. \\
C & Special-function construction & Generic genus-two basis and horizon letter. \\
D & Algebraic degeneration (or supplement) & Negative-radius genus-one limit and detailed residue/
nome machinery. \\
\bottomrule
\end{longtable}

This structure leaves the science intact while making the main line legible.

\section{Figures, tables and visual presentation}

\subsection{Global visual assessment}

The PDF is technically clean: there are no missing plots, clipped equations or visibly corrupted
characters.  The body font is readable at full-page size and the running header is consistent.
The main concern is proportionality and final-size legibility.  IOP advises authors to use clear,
high-quality figures, text around 8--12 points at final size, and non-colour distinctions.  Several
plots meet the data requirement but not yet the journal-design requirement.

\subsection{Figure-by-figure recommendations}

\begin{longtable}{@{}p{0.10\textwidth}p{0.22\textwidth}p{0.58\textwidth}@{}}
\toprule
Item & Current role & Editorial action \\
\midrule
\endhead
Fig. 1 & Three indicatrix geometries & Redraw as vector schematic with larger labels.  Remove or
visually demote Thakurta--Kerr.  Use panel labels and a single legend. \\
Fig. 2 & FLRW freezing/redshift & Combine with figure 3.  Remove the long in-plot title; put the
interpretation in the caption. \\
Fig. 3 & FLRW perturbation check & Relabel as a restricted local perturbation test, not a proof of
global minimisation. \\
Fig. 4 & Vaidya indicatrix & Keep if it directly introduces the radial wind; harmonise fonts and
colours with figure 1. \\
Fig. 5 & Branch depth/penetration & Keep; make exact/numerical/frozen status explicit in caption. \\
Fig. 6 & Timing/trajectory behaviour & Keep the strongest panel; ensure the formal branch is not
visually equated with physical evaporation. \\
Fig. 7 & No-inversion diagnostic & Move to appendix/supplement.  Its caption is mostly a disclaimer
and the main result is negative. \\
Fig. 8 & Bounce through periapsis & Keep, but place before the Conclusions and reduce plot title. \\
Fig. 9 & Arrival-epoch timing & Keep or combine with figure 8; state that the blue branch is formal. \\
Fig. 10 & Off-shell residual scaling & Definitely keep.  Add fit interval, fitted slopes with
uncertainty, norm, tolerances and sub-arc.  Correct any stale equation reference. \\
Fig. A1 & Stationary-limit residual & Keep in reproducibility appendix at smaller size. \\
Fig. A2 & Branch perturbations & Move to supplement or re-caption as a local family check. \\
\bottomrule
\end{longtable}

Use line styles and markers in addition to colour so that physical/static/formal branches remain
distinguishable in greyscale.  Export line art as vector PDF/EPS where possible.  Avoid titles such
as ``gap crosses 0 ONLY inside $m<0$'' inside the graphics; that sentence belongs in prose.

\subsection{Float order}

No figure should interrupt the Conclusions.  Insert a float barrier or explicit clear-page before
section 5 after placing figures 8--10.  A better solution is to integrate figures 8 and 9 into the
phenomenology subsections and place figure 10 at the end of the first-order section.  Then begin the
Conclusions on a fresh page if necessary.

\subsection{Tables}

Table 1 should become a genuine Paper-I notation table.  Table A1 can be shortened.  Table A2 must
be rebuilt because the final rows visibly collide.  Suggested columns are: Result; equation/
figure; script; archived hash; verification quantity; achieved value.  Do not put a prose category
header inside a row with data.

\section{Notation and terminology consistency}

The mathematical notation is mostly stable, but the number of clocks and nearby uses of the same
symbols create an avoidable barrier.  A compact convention table should distinguish:

\begin{tabularx}{\textwidth}{@{}p{0.14\textwidth}X@{}}
\toprule
Symbol & Required unambiguous use \\
\midrule
$t$ & Stationary/frozen coordinate time only; not a canonical clock after $m=m(v)$. \\
$\eta$ & FLRW conformal time; do not reuse descriptively for a generic slow clock. \\
$\tau$ & Proper time along the rail. \\
$v$ & Advanced null coordinate and native ingoing-Vaidya evolution parameter. \\
$u$ & Four-velocity in most of the paper; avoid using $u$ for retarded time without an explicit
font/context warning. \\
$s$ & Abstract selected evolution clock or slow parameter. \\
$\lambda$ & Elapsed clock anchored at the free-arrival endpoint; always state orientation. \\
$J$ & Conserved axial Pontryagin costate. \\
$L_{\rm mech}$ & Mechanical angular momentum, generally not conserved on the controlled rail. \\
$J_{\rm deg}$ & Algebraic genus-degeneration value in Vaidya; not automatically a physical
separatrix. \\
$J_c$ & Use only for a physically defined critical/capture threshold, or define explicitly if it
is deliberately identified with $J_{\rm deg}$ in a restricted context. \\
\bottomrule
\end{tabularx}

The manuscript sometimes uses ``arrival branch'', ``coordinate-time branch'', ``advanced-time
branch'' and ``native branch'' close together.  Define branch names once per geometry.  For Vaidya,
prefer ``advanced-time ($v$) branch'' and ``proper-time ($\tau$) branch''; call $t$ only the frozen
Schwarzschild comparison clock.

Use ``extremal'' for a PMP solution, ``locally minimizing within the tested family'' for a numerical
perturbation check, and ``global minimizer'' only when the HJB/cut-locus conditions are supplied.
Use ``closed representation'' when numerical periods/theta evaluation remain necessary, and
reserve ``elementary/elliptic closed form'' for the stronger cases.

The manuscript alternates ``non-stationary'' and ``non-autonomous''.  They are related but not
synonymous: the former describes the geometry, the latter the reduced control system.  Maintain
that distinction intentionally.

\section{Bibliography, data and research-integrity audit}

\subsection{Reference-list curation}

The source currently calls \verb|\nocite{*}| and two bibliography databases.  This produces 69
printed references, while only 40 keys are explicitly cited in the Paper-I source.  The unused set
includes predecessor projects, Kerr geodesics, Thakurta/McVittie cosmological black holes,
classical algebraic-geometry books and recent polylogarithm papers not invoked in Paper I.  This is
not merely cosmetic: it makes the article look like an incompletely separated master manuscript.

Remove \verb|\nocite{*}|.  Compile, inspect all undefined citations, and then add only works that
are discussed in the text.  Replace irrelevant carry-over references with the missing direct
brachistochrone literature.  The final list may be shorter while being substantively stronger.

Every entry should be checked against the version of record.  Prefer DOI links; include titles if
the selected IOP style permits or if they materially help a specialised bibliography.  Software
references should identify exact versions and persistent archives.  If Paper II is not publicly
available at submission, follow the journal's rule for unpublished material rather than treating
it as an ordinary retrievable paper.

\subsection{Data and software}

The Zenodo DOI is an important strength.  Verify four things immediately before submission:

\begin{enumerate}
  \item the DOI resolves publicly without authentication;
  \item the cited record/version contains every script named in Table A2;
  \item the hashes printed in the PDF match that archive, not the development branch;
  \item the environment records Python/Sage/Mathematica and package versions needed for the
  period-matrix and theta calculations.
\end{enumerate}

The GitHub link should be supplementary to, not a replacement for, the immutable archive.

\subsection{Generative-AI disclosure}

The disclosure is broadly consistent with current IOP policy: it names systems and versions,
describes uses, states critical review and assigns responsibility to the author.  Keep it.  The
author should also retain private records sufficient to explain which reference metadata and
symbolic outputs were independently verified, because IOP places responsibility for accuracy on
the author and does not accept unverified AI-generated reference lists.

No authorship or competing-interest concern is apparent from the reviewed material.  The single-
author contribution statement is implicit; a short CRediT statement is optional rather than
necessary.

\section{Major-revision checklist}

The following items are ordered by editorial impact.  Items 1--10 should be completed before
submission; the remaining items are strongly recommended.

\begin{enumerate}[label=\textbf{M\arabic*.}]
  \item \textbf{Repair the novelty paragraph.}  Remove the absolute ``have not been addressed''
  claim, add the direct 1997--2002 relativistic-brachistochrone literature, and state the controlled
  non-autonomous/Kodama advance precisely.
  \item \textbf{Promote Theorem I.5.}  Put its statement and interpretation in the main text;
  retain the derivation in an appendix.  Restore result numbering in reading order.
  \item \textbf{Shorten section 4.4.}  Keep the complete response, universal source, representation
  hierarchy and true-flow validation in the body.  Move detailed divisor, degeneration, nome and
  minimal-basis material out of the main narrative.
  \item \textbf{Make optimality language uniform.}  Correct ``minimizes'' in figures 3 and A2 and
  qualify ``well-posed''/``global'' everywhere by the assumptions already stated in section 2.
  \item \textbf{Complete the Paper-I/Paper-II split.}  Remove most Thakurta--Kerr notation and
  detailed ergosphere/conformal-momentum discussion from Paper I; retain only a concise ladder and
  forward reference.
  \item \textbf{Regenerate the bibliography.}  Delete \verb|\nocite{*}|, cite only relevant work,
  add missing direct predecessors and validate all metadata.
  \item \textbf{Rebuild figure hierarchy.}  Combine FLRW plots, move figure 7 to supplement, place
  figures 8--10 before the Conclusions and make all labels legible at final size.
  \item \textbf{Repair Table A2.}  Eliminate overlapping content, separate scripts/hashes/tests and
  tie every row to the archived release.
  \item \textbf{Clarify endpoint terminology.}  Use fixed launch event plus fixed spatial target at
  free arrival clock consistently; distinguish this from a two-event fixed-spacetime-endpoint BVP.
  \item \textbf{Calibrate ``closed form''.}  Use ``closed special-function representation'' where
  period construction and archived numerical routines are still required; keep conjectural
  irreducibility explicitly labelled.
  \item \textbf{Re-centre the abstract.}  Reduce terminology, identify the physical Vaidya result,
  and avoid implying an unconditional global minimisation theorem.
  \item \textbf{Explain the rail operationally in the introduction.}  State what actuation is
  allowed, what charge is fixed and which cost is minimised.
  \item \textbf{Separate result levels in Vaidya.}  Mark exact, frozen, numerical, adiabatic and open
  outgoing statements with a consistent visual/verbal scheme.
  \item \textbf{Standardise $J_c$ versus $J_{\rm deg}$.}  Do not call a negative-radius algebraic
  degeneration a physical separatrix.
  \item \textbf{Recompile and compare.}  Resolve the stale equation reference seen in figure 10's
  supplied-PDF caption and check every cross-reference after reordering.
\end{enumerate}

\section{Minor-revision checklist}

\begin{enumerate}[label=\textbf{m\arabic*.}]
  \item Choose one title and a stable five-to-seven-keyword set.
  \item Use one spelling convention for ``non-stationary'' and consistent hyphenation for
  ``controlled-rail'', ``free-arrival'' and ``first-order''.
  \item Define CKV in words at first use in the body even if the abstract spells it out.
  \item Reduce parenthetical caveats by moving them into short remarks or a limitations subsection.
  \item Break long paragraphs in section 4.4; several contain multiple logical levels and formulas.
  \item State whether $\dot m$ is dimensionful and how $\varepsilon=M\dot m/m$ is used in every
  residual plot.
  \item Give fit uncertainty and fitting interval for slopes in figure 10, not only ``approximately
  2.00''.
  \item Name the residual norm and integration sub-arc in each validation caption.
  \item Use panel labels (a), (b) instead of relying on vertical position.
  \item Replace all-caps emphasis inside plots with ordinary sentence case.
  \item Ensure line styles, not colour alone, distinguish branches.
  \item Export line art in vector form and confirm 8--12 pt final text.
  \item Remove unused abbreviations and symbols from Table 1.
  \item Give $W$, $s$ and the terminal manifold in one protocol box.
  \item Make the orientation of $\lambda$ explicit in the first equation where it appears, not only
  in Appendix C.
  \item Avoid using $u$ both for four-velocity and retarded time in the same local discussion.
  \item Distinguish ``horizon'', ``apparent horizon'', ``trapping horizon'' and algebraic branch
  point according to the Vaidya assumptions.
  \item Keep ``formal negative-rate continuation'' in every relevant legend/caption.
  \item Shorten captions: describe what the reader should learn, not the full derivation and all
  limitations.
  \item Place computational provenance (script name/hash) in Appendix A rather than every main
  caption, unless essential.
  \item Make equation numbering stable after moving Theorem I.5; avoid hard-coded equation numbers
  in plot artwork.
  \item Add a single paragraph explaining why the genus-two structure is surprising physically,
  rather than only contrasting it with geodesic ellipticity.
  \item State once that additive constants/reference lengths in logarithms cancel in assembled
  observables; do not repeat.
  \item Move the rank-five numerical evidence to the archive or supplement and label its data set.
  \item Verify that all manuscript software names use their official capitalisation.
  \item Shorten the Paper-II discussion in the abstract, introduction, Table 1, figure 1, section 2,
  Appendix C and Conclusions to avoid repetition.
  \item Consider a dedicated ``Limitations'' subsection before Conclusions rather than distributing
  caveats across captions.
  \item Ensure the data-availability DOI points to the version corresponding to the final PDF.
  \item Retain the AI disclosure but separate it typographically from software acknowledgements.
  \item Run a final English copy-edit focused on article use, sentence length and repeated
  ``genuinely'', ``cleanly'', ``exactly'' and ``universal''.
\end{enumerate}

\section{What should be left substantially unchanged}

Revision should not flatten the paper into a generic optimal-control article.  The following
elements are distinctive and already effective:

\begin{itemize}
  \item the fixed charge $-u\cdot W=\hat E$ as the central controlled-rail definition;
  \item the Killing/CKV/Kodama selector hierarchy, presented as a declared geometric convention;
  \item the explicit statement that the rail force may perform work against $W$ and is not free
  fall;
  \item the unique-maximising-direction PMP reduction on a compact strictly convex indicatrix;
  \item the distinction between endpoint transversality $H(s_f)=0$ and the conserved extended
  Hamiltonian;
  \item the separation of $J=p_\varphi$ from mechanical angular momentum;
  \item the conditional HJB verification language and cut/Maxwell qualification;
  \item FLRW as the centred, branch-degenerate sanity check;
  \item the unnormalised Kodama selector in ingoing Vaidya;
  \item the distinction between the native advanced-time branch and the proper-time branch;
  \item the formal/nonphysical status of the negative-rate ingoing continuation;
  \item the complete on-shell plus off-shell first-order correction and its true-flow residual test;
  \item the universal Vaidya boundary source $S_{\mathcal D}=[r p_r]-\lambda$;
  \item the ``Proved / Conditional and numerical / Outlook'' result-tiering;
  \item the immutable data/code archive and transparent AI-use disclosure.
\end{itemize}

These are the elements a revised article should make easier, not harder, to see.

\section{Model editorial text}\label{sec:modeltext}

The following examples illustrate claim calibration.  They are editorial proposals, not
scientific rewrites, and should be checked against the final notation.

\subsection{Suggested abstract}

\begin{quote}
We formulate a time-optimal controlled-worldline problem for spacetimes without a timelike
Killing symmetry.  An ideal rail maintains the invariant $-u\cdot W=\hat E$, while the control is
the direction on the resulting velocity indicatrix and the arrival clock is free.  The selector
$W$ is chosen from a Killing--conformal-Killing--Kodama hierarchy, reducing to the Kodama vector in
dynamical spherical symmetry.  On compact domains where the selector is timelike and the
indicatrix is nondegenerate, we derive the normal Pontryagin extremals, establish existence under
the stated coercivity assumptions, and give a time-dependent Hamilton--Jacobi criterion for global
verification.  Spatially flat FLRW provides a branch-degenerate base case: the admissible
indicatrix is centred and all considered clock functionals have the same straight comoving spatial
path.  In ingoing Vaidya, radial mass flow shifts the indicatrix, separates the advanced-time and
proper-time branches, and makes penetration depend on the arrival epoch.  For slowly varying mass
we derive the complete first-order response, including the interior off-shell costate term.  Its
Vaidya source reduces to a boundary expression, and the resulting Abelian/iterated-integral
representation is verified against the full non-autonomous flow with an $O(\varepsilon^2)$
residual.  The outgoing-Vaidya boundary-value problem is left open.
\end{quote}

This version is approximately the same length as the current abstract but removes Paper-II detail,
states the domain and makes the physical result visible.

\subsection{Suggested end of the novelty paragraph}

\begin{quote}
Relativistic brachistochrones and arrival-time principles have a substantial stationary and
geometric variational literature.  The difficulty addressed here is different: the prescribed
charge is actively maintained while the background geometry evolves, so the reduced control
problem is non-autonomous and its ordinary Hamiltonian is not an interior first integral.  In
spherical symmetry the Kodama vector supplies a canonical selector for this controlled problem.
\end{quote}

\subsection{Suggested claim-status box}

\begin{quote}
\textbf{Scope of the result.}  The PMP equations characterise normal extremals on the regular
timelike-selector domain.  Global minimisation follows only where a single-valued HJB verification
function or an explicit comparison across competing extremals is available.  The Vaidya
negative-rate curves shown below are formal continuations in the ingoing chart and are not a model
of physical evaporation.
\end{quote}

Placed near the end of section 2, this box would allow later captions to remain short.

\subsection{Suggested conclusion core}

\begin{quote}
The controlled-rail constraint gives a well-defined non-autonomous time-optimal problem on compact
domains where the selected field is timelike and the velocity indicatrix is nondegenerate.  FLRW
is the centred base case, whereas the Kodama flow in Vaidya shifts the indicatrix radially and makes
the optimal branches sensitive to the mass history and arrival epoch.  The complete first-order
adiabatic response requires both the drift of the frozen family and the off-shell costate
displacement.  For Vaidya the latter has a universal boundary source, and the assembled correction
agrees with the full non-autonomous flow up to quadratic residuals.  Global HJB certification,
turning-point uniformisation and the physical outgoing-Vaidya control problem remain open.
\end{quote}

\section{A practical revision sequence}

The work can be revised efficiently without reopening the science in every section.

\begin{enumerate}
  \item \textbf{Freeze the claim hierarchy.}  Write a one-page internal sheet listing each main
  result as proved, conditionally verified, numerical, conjectural or open.  Use that sheet to
  audit the abstract, captions and conclusion.
  \item \textbf{Repair literature and introduction first.}  This determines the exact novelty
  claim and prevents later sections from being framed too broadly.
  \item \textbf{Move Theorem I.5 next.}  Once the main theorem is in the body, reorganise section
  4.4 around it and move technical derivations outward.
  \item \textbf{Complete the paper split.}  Delete or condense Paper-II-specific notation and
  passages before editing prose; otherwise the same material will be polished twice.
  \item \textbf{Rebuild figures and tables.}  Decide which plots are main results, appendix checks
  or supplementary diagnostics.  Then fix float order and captions.
  \item \textbf{Regenerate references and cross-references.}  Remove \verb|\nocite{*}|, compile
  twice, inspect warnings and compare the PDF against the archived code map.
  \item \textbf{Perform a final broad-reader pass.}  Read only title, abstract, introduction,
  first sentence of each section, captions and conclusion.  That reduced document should tell a
  complete, accurately qualified story.
\end{enumerate}

This sequence is preferable to line-editing the current 38 pages from top to bottom because the
main defects are structural.  Once the structure is fixed, the sentence-level edit will be much
smaller.

\section{Final recommendation}

The manuscript should be revised, not abandoned.  The central controlled-rail idea is distinctive;
the use of Kodama flow in a dynamical spherical spacetime is well matched to CQG; the endpoint and
off-shell issues are treated far more carefully than in earlier development notes; and the
true-flow residual provides a concrete validation target.  There is enough substance for a strong
Paper I.

I would not submit the supplied PDF unchanged.  Its present framing exposes the paper to avoidable
criticism on novelty, global optimality language and physical interpretation, while its length and
special-function density conceal the result that is most likely to interest the journal.  The
unused 69-entry bibliography, Paper-II notation in Table 1, out-of-order theorem numbering,
caption overclaims, broken Table A2 and figure-fragmented Conclusions reinforce the impression of a
development snapshot rather than a final journal article.

My simulated editorial recommendation is therefore:

\begin{quote}
\textbf{Return for major editorial revision, with encouragement to resubmit.}  A revised version
should be suitable for external specialist review if it (i) positions the novelty against the
direct relativistic-brachistochrone literature, (ii) centres the complete extended-Hamiltonian
Vaidya response in the main narrative, (iii) moves most higher-genus construction detail out of
the body, and (iv) aligns every claim, caption, figure and reference with the paper's carefully
stated domain and optimality hierarchy.
\end{quote}

If those four conditions are met, the remaining issues are ordinary copy-editing and production
work.  The author should then proceed with the physics rather than reopen already settled algebra
unless a mathematical referee identifies a specific defect.

\bigskip
\section*{Sources consulted for the editorial assessment}
\addcontentsline{toc}{section}{Sources consulted for the editorial assessment}

\begin{itemize}
  \item Classical and Quantum Gravity author guidelines and scope:
  \url{https://publishingsupport.iopscience.iop.org/journals/classical-and-quantum-gravity/}
  \item IOP generative-AI policy:
  \url{https://publishingsupport.iopscience.iop.org/questions/generative-ai-tools/}
  \item F. Giannoni, P. Piccione and D. V. Tausk, \emph{Morse theory for the travel time
  brachistochrones in stationary spacetimes}, DCDS 8 (2002) 697--724,
  \url{https://doi.org/10.3934/dcds.2002.8.697}
  \item F. Giannoni and P. Piccione, \emph{The arrival time brachistochrones in general
  relativity}, Journal of Geometric Analysis 12 (2002) 375--423,
  \url{https://doi.org/10.1007/BF02922047}
  \item E. Caponio and M. A. Javaloyes, \emph{On the formulations of the Fermat principle in
  general relativity and beyond}, arXiv:2605.01532,
  \url{https://arxiv.org/abs/2605.01532}
  \item \emph{The brachistochrone problem for a constant velocity traveler in static and
  spherically symmetric spacetimes}, arXiv:2607.29265,
  \url{https://arxiv.org/abs/2607.29265}
\end{itemize}

\vfill
\begin{center}
  \small End of editorial report.
\end{center}

\end{document}
